\chapter{From Predecessor Prototypes to the Research PoC}
\label{chap:legacyAnalysis}

Twin2MultiCloud did not start as a greenfield product. It combines and narrows
ideas from two independently developed prototypes: a browser-based cost
optimizer for a Five-layer Digital Twin architecture and an AWS-oriented cloud
deployer. Their code and assumptions were useful starting points, but their
outputs were not connected and neither prototype could provide the evidence
required by the research questions. This chapter derives the transformation
decisions that led to the Six-layer proof of concept (PoC). The resulting
architecture is described separately in Chapter~\ref{chap:systemArchitecture}.

\section{Predecessor Cost Optimizer}
\label{sec:legacy-optimizer}

The predecessor optimizer was a client-side JavaScript application. It read
scenario values from browser controls, calculated costs for the responsibilities
of a Five-layer model, and rendered the cheapest assignment into the document
object model. This implementation established the central idea of selecting
cloud providers per architectural responsibility and remains valuable as the
offline \emph{Five-layer-v1} baseline.

Four properties prevented the prototype from directly answering the current
research questions:

\begin{enumerate}
    \item Its result was a presentation object rather than a versioned
    machine-readable architecture and deployment contract.
    \item Cost calculation and user-interface state were coupled, which made
    deterministic reproduction outside a browser difficult.
    \item Provider price values and assumptions were not represented as a
    separately attributable evidence set.
    \item Monetary ranking did not prove that the selected provider services
    formed a functionally complete, deployable graph.
\end{enumerate}

The baseline therefore remains an offline comparison method; it is not an
alternative runtime profile. The active Optimizer implements only
cost-minimization for \texttt{six-layer-eventing@1}. A small internal strategy
boundary preserves the software pattern needed for a later objective, but no
latency, sustainability, weighted-scoring, or public strategy-selection feature
is shipped in the PoC.

\section{Predecessor Cloud Deployer}
\label{sec:legacy-deployer}

The predecessor Deployer provisioned an AWS-oriented Five-layer architecture
through provider SDK operations and configuration files. It demonstrated that
the layered model could be materialized, but it did not consume an optimizer
result and did not implement equivalent Azure and GCP deployments. Shared
mutable configuration and imperative resource creation also made it difficult
to associate a deployment with one immutable calculation and to prove cleanup.

The relevant gaps were:

\begin{enumerate}
    \item no common contract linked workload, architecture selection, exact
    provider resources, verification, and destruction;
    \item layer names suggested provider variability, while the implemented
    resource bundles did not establish three-provider functional equivalence;
    \item secrets and non-secret project configuration did not have a clear
    ownership and transport boundary;
    \item partial failures could leave resources whose relationship to the
    originating request was not represented as durable evidence; and
    \item the executable unit was an editable project directory rather than a
    bounded, generated deployment-operation package.
\end{enumerate}

The PoC consequently uses Terraform for declarative provider resources and a
short-lived operation package generated by the Management API. The internal
workspace is still called a project in parts of the implementation, but it is
not a user-managed project abstraction: users create and transfer typed Twins,
whereas the Deployer only stages one generated package for a specific Twin and
operation.

\section{Research-Driven Gap Analysis}
\label{sec:legacy-gap-analysis}

The predecessor gaps are relevant only insofar as closing them creates evidence
for the research questions. Table~\ref{tab:gap-rq-map} records this trace rather
than defining a general product backlog.

\begin{table}[htbp]
\centering
\caption{Predecessor gaps and the bounded response in the research PoC.}
\label{tab:gap-rq-map}
\begin{tabular}{p{3.2cm}p{4.7cm}p{4.7cm}}
\hline
\textbf{Gap} & \textbf{Required evidence} & \textbf{PoC response} \\
\hline
Disconnected calculation and deployment & Reproduce how one input becomes a
deployed and destroyed graph (RQ1) & Digest-linked workload, calculation, RTA,
RDS, manifest, operation, verification, and cleanup evidence \\

Nominal rather than proven provider equivalence & Compare complete provider
implementations without assuming identical services (RQ2) & Curated component
bundles, capability profiles, edge contracts, negative candidates, and runtime
probes \\

Cost winner without deployability proof & Compare only admissible provider
assignments (RQ2 and RQ3) & Functional resolution precedes exact cost scoring;
incomplete candidates are rejected \\

No comparable three-provider experiment & Measure local and directed
multi-cloud cases under one workload (RQ3.1) & Three provider-local baselines
and six directed provider-pair scenarios \\

Event transport hidden inside other layers & Assess the architectural and cost
effect of explicit eventing (RQ3.2) & Standalone Six-layer contract plus an
offline Five-layer-v1 baseline \\
\hline
\end{tabular}
\end{table}

This mapping also determines what the implementation does \emph{not} need.
Production authentication, multi-tenancy, account creation, generic architecture
plugins, multiple active objectives, arbitrary executable project archives,
in-place infrastructure updates, dashboard administration, and exhaustive layer
permutations would add engineering surface without directly improving the
planned evidence.

\section{Transformation Method}
\label{sec:refactoring-methodology}

The transformation followed an evidence-first sequence. A contract was defined
before its implementation boundary; provider-specific behavior was admitted
only when it could satisfy that contract; deterministic tests were added before
live evaluation; and cross-service outputs were made immutable and
digest-linked. This sequence is more important to the thesis than the number of
implemented endpoints or design patterns.

The refactoring also used small extension boundaries where variation is
structurally plausible. Examples include provider adapters, price fetchers,
resource calculators, cost scoring, and validated user-function slots. In the
active PoC, however, an extension boundary does not imply multiple shipped
variants. One cost strategy and one architecture profile are sufficient to
demonstrate the pattern without claiming a generic framework.

\subsection{Optimizer Transformation}
\label{sec:refactoring-optimizer}

The browser-bound calculations were moved behind a typed Python API. Pricing is
represented by frozen, dated provider snapshots with source attribution and
digests. The calculation pipeline first enumerates deterministic assignments,
then applies architecture and capability constraints, materializes the exact
deployment specification, and only then ranks admissible candidates by monthly
cost. Its output is consequently suitable for both cost comparison and
deployment traceability.

The original Five-layer-v1 method remains executable only as an offline
baseline. The Six-layer profile is independently versioned and does not inherit
from another executable profile. Its explicit Eventing responsibility can
therefore be compared with the historical baseline without presenting two
selectable deployment architectures.

\subsection{Deployer Transformation}
\label{sec:refactoring-deployer}

Provider behavior was separated from orchestration and the executable
infrastructure was expressed in Terraform. AWS, Azure, and GCP implement the
fixed profile through curated provider service bundles rather than through a
claim that similarly named services are equivalent. The Deployer validates the
resolved graph and operation package, executes Apply or Destroy, and emits
verification and cleanup evidence.

The public Deployer boundary consists of typed configuration and bounded
function validation plus the generated operation package used by Management.
Its internal working directory is transient execution state rather than a
second user-facing model next to the Twin.

\subsection{Management and User Workflow}
\label{sec:refactoring-management}

The Management API and Flutter client connect the two former prototypes. A
Twin owns typed configuration and bounded function sources. A user may export,
import, or duplicate a Twin through an allowlisted \texttt{.twin.zip}, but the
archive never contains cloud credentials or deployment state. CloudConnections
are encrypted, profile-owned records selected for the providers required by a
resolved graph.

One local owner profile preserves this ownership relationship. Interactive
Google, Microsoft, or university-SAML application login, JWT issuance, roles,
and multi-tenant sessions were removed because they do not answer the research
questions. The compiled but unrouted login page is only an adapter boundary for
future work. Provider-side IAM, Entra, IAP, OIDC, or SAML remains a separate
deployment concern where a resolved cloud graph requires it.

Deploy and Destroy remain explicit, durable operations because they mutate
external resources and can incur cost. Progress is persisted and streamed with
reconnect support, while retries are protected by idempotency. A deployed Twin
is immutable; changes produce a new draft and calculation rather than an
in-place Terraform update or migration.

\section{Resulting Evaluation Boundary}
\label{sec:transformation-evaluation-boundary}

The final transformation target is a reproducible Small workload evaluated in
three provider-local scenarios and all six directed AWS/Azure/GCP provider
pairs. These nine scenarios cover provider locality and every cross-provider
direction without enumerating every possible layer assignment. Structural
constraints, such as the required co-location of hot storage and visualization,
are tested offline as negative candidates instead of being forced into invalid
live deployments.

Live execution is a separate empirical phase. Every scenario requires an
approved numerical budget cap, one active deployment at a time, bounded
duration, explicit Apply and Destroy confirmation, functional verification,
and a post-Destroy provider inventory. Until those records exist, the
repository may claim deterministic contract and candidate evidence, but not
successful multi-cloud operation or observed cost savings.

The refactoring and its decisions were AI-assisted. Repository inspection,
implementation, tests, and documentation used AI tooling, while scope choices,
source assessment, research interpretation, credentials, cost approval, and
live-cloud authorization remain human responsibilities. AI-generated output is
treated as an engineering proposal to validate, not as empirical evidence or a
scientific source.
